\section{Results}
\label{sec:results}

This section reports the experimental setup and the diagnostics
produced by \rolloutscript{}.
We use identical-twin experiments throughout: the truth and the
assimilation share the same forward model class, but the truth is
driven by a parameter trajectory drawn from a separate generator
(Section~\ref{subsec:da-prior}).

\subsection{Experimental setup}
Unless otherwise stated, the configuration follows
\code{scripts/config.py}:
\begin{itemize}
  \item \emph{Horizon and windows.} $\Nw = 6$ assimilation windows of
        length $\windowLen = 300\,\mathrm{s}$ each, for a total
        horizon of $30$ minutes.  Each window represents the
        parameter trajectory on $\Nt = \windowLen / 60 = 5$
        time points, sharing boundary points with neighbors.
  \item \emph{Solver.} Both truth and assimilation use the
        \pylbm{} solver on a $50 \times 40 \times 16$ grid covering
        the staggered cube array of Xie and
        Castro~\cite{xie2008castro}.
  \item \emph{Smoother.} \ESMDA{} with $\Nsteps = 3$ tempering
        steps, $\alphaInfl_i = 3$, and ensemble size $\Ne = 32$.
        \code{pin\_initial\_time\_point} is disabled in window $0$
        and enabled from window $1$ onward.
  \item \emph{Observations.} Five probe locations sampling the
        three velocity components $(u, v, w)$ near the ground, with
        observation-noise standard deviation $\sigma_d = 0.25$ and
        time-mean aggregation over fixed-length intervals.
  \item \emph{External prior.} $\xext$ and $\Sext$ for the inflow
        angle (mean $0^\circ$, std $10^\circ$) and the inflow
        speed (mean $5\,\mathrm{m\,s^{-1}}$, std $0.5\,\mathrm{m\,s^{-1}}$,
        clipped at $0.1\,\mathrm{m\,s^{-1}}$).
  \item \emph{Prior model.} The default critically-damped AR(2)
        relaxation of Section~\ref{subsec:da-prior}, with
        correlation length $\ell_{\mathrm{prior}} =
        200\,\mathrm{s}$ for the assimilation prior and
        $\ell_{\mathrm{truth}} = 200\,\mathrm{s}$ for the truth.
        Different correlation lengths are easily configured for
        sensitivity studies.
\end{itemize}

\subsection{Diagnostics}

\paragraph{Time-varying parameter recovery.}
The headline diagnostic is the per-window prior, posterior, and
truth trajectories of each parameter, plotted against absolute time
by the routine \plothelper{}.  The
figure overlays:
\begin{itemize}
  \item the truth concatenated across windows (solid blue line);
  \item the per-window prior ensemble -- cold-start for window $0$,
        AR(2)-extrapolated for $w > 0$ -- as dashed green lines
        (individual members) and a thick dashed green ensemble mean;
  \item the per-window posterior ensemble in solid orange
        (individual members) with a thick orange ensemble mean.
\end{itemize}
Window boundaries are marked with vertical dashed gray lines.  The
plot is saved to \code{time\_varying\_parameters\_rollout.png} in the
run-specific output directory and is the primary visual evidence
that the smoother tracks the truth across the full horizon.

\paragraph{Expected behavior.}
The cold-start prior of window $0$ has the full external spread
$\Sext$ and is centered on $\xext$; the smoother pulls the mean
toward the truth and contracts the spread.  From window $1$ onward
the prior at $t = 0$ inherits the previous-window posterior mean
through equation~\eqref{eq:ar2-relax} and relaxes back to $\xext$ on
the time scale $\ell_{\mathrm{prior}}$, while \ESMDA{} keeps the
seam value pinned during the update.  As windows accumulate, the
posterior trajectory should track the truth with a window-to-window
continuity that is visible to the eye.

\paragraph{State diagnostics.}
Two additional diagnostics check that the recovered parameters
translate into a faithful flow:
\begin{itemize}
  \item \code{state\_init\_and\_terminal.png} -- horizontal slices of
        the truth and the posterior ensemble mean at the initial and
        terminal times of the global trajectory, with probe
        locations overlaid;
  \item \code{rollout\_animation.mp4} -- a side-by-side animation of
        the truth and the ensemble-mean state across the full
        $\Nw\windowLen$ horizon, useful for inspecting the
        behavior at window transitions.
\end{itemize}

\paragraph{Saved data.}
For reproducibility, the rollout writes five NetCDF archives at the
end of the run: \code{true\_state.nc}, \code{esmda\_state.nc},
\code{true\_params.nc}, \code{posterior\_params.nc}, and
\code{prior\_params.nc}, each on absolute time.  These are the
inputs to all diagnostic plots and provide a stable interface for
post-hoc analyses (e.g.\ RMSE versus the truth, posterior
calibration, spectra of the recovered parameters).

\subsection{Status}
At the time of writing, this report serves as a methodological
document; concrete figures and quantitative comparisons across prior
models, correlation lengths, and solver pairings will be filled in
during the campaign that produces the academic article.  The
ablations envisaged for that campaign are listed in
Section~\ref{sec:conclusion}.
